% =====================================================================
%  Seksjon: Skaft-/generatormoment-eksitasjon av tårnmodene i begge
%           modeller (forenklet WindTurbineTower + OpenFAST-FMU), inkl.
%           manipulering av ROSCO open-loop-fila, og SS-vs-fore-aft-
%           selektiviteten.
%  Kilder:  src/tops_openfast/dyn_models/windturbine_tower.py
%           casestudies/dyn_sim/sweep_resonance.py (FMU, OL-table)
%           casestudies/dyn_sim/_contrast_ss_fa.py
%           test1002/ControlData/ROSCO.IEA15MW.IN  (OL_Mode, Ind_GenTq)
%  Preamble: amsmath, amssymb, graphicx, booktabs, siunitx, hyperref
% =====================================================================

\section{Shaft (Generator-Torque) Excitation of the Tower Modes}
\label{sec:shaft-excitation}

This project studies electro-mechanical interactions in offshore wind turbines
coupled to a grid --- here one connected to the LEOGO islanded platform network.
One important such interaction is whether an \emph{electrical} disturbance can
excite a \emph{mechanical} tower mode. The physical link is the generator/air-gap
electromagnetic torque acting through the shaft: fluctuations in the delivered
power \(P_e\) become fluctuations in \(T_e = P_e/(\omega_{e,\mathrm{filt}}\eta)\),
whose nacelle reaction torque bends the tower side to side. This section collects
how that shaft-torque excitation is realised in the two models --- the reduced
\texttt{WindTurbineTower} of \S\ref{sec:simplified-tower-model} and the
OpenFAST-FMU of \S\ref{sec:openfast-excitation} --- and shows that the torque
path selectively reaches the side-to-side (SS) mode but not the fore-aft (FA)
mode.

\subsection{Reduced model: an open electrical-to-mechanical path}
\label{subsec:se-reduced}

In \texttt{WindTurbineTower} the side-to-side oscillator \eqref{eq:stm-ss} is
driven directly by \(T_e\), and \(T_e\) is computed from the power the converter
actually delivers to the grid. The path
\begin{equation}
  \text{grid disturbance}\;\longrightarrow\;P_e
  \;\longrightarrow\;T_e
  \;\longrightarrow\;\text{nacelle reaction moment}
  \;\longrightarrow\;q_{\mathrm{ss}}
  \label{eq:se-path}
\end{equation}
is therefore \emph{open by construction}: no controller intercepts it, unlike in
the FMU. A sustained sinusoidal grid load at the modal frequency thus feeds the
SS mode continuously and drives it into resonance, and a swept load traces out
its resonance curve. Both are demonstrated quantitatively in
\S\ref{sec:new-tests}; here the point is only that the reduced model provides the
genuinely closed electrical\(\to\)mechanical demonstration that the FMU cannot,
because in the FMU ROSCO owns the torque.

\subsection{OpenFAST model: injecting torque through the ROSCO open-loop file}
\label{subsec:se-openfast}

In the OpenFAST build used here ROSCO (\texttt{VSContrl}=5) computes the
generator torque internally from the measured generator speed and rejects every
external torque or power command passed across the FMU interface
(\texttt{GenSpdOrTrq}, \texttt{ElecPwrCom}); this is the negative result
documented in \S\ref{subsec:of-electrical}. Consequently the network cannot reach
the OpenFAST \emph{structure} through the normal electrical inputs, and the only
channel that does reach the structural solver is ROSCO's own \emph{open-loop
control} table.

\paragraph{What was changed in the ROSCO input.}
ROSCO reads its input files at runtime, so the torque channel was opened without
rebuilding the FMU, by editing \texttt{test1002/ControlData/ROSCO.IEA15MW.IN}:
\begin{itemize}
  \item \texttt{OL\_Mode}: \(0\to1\) --- enable open-loop control versus time.
  \item \texttt{Ind\_GenTq}: \(0\to2\) --- take the generator torque from column~2
        (in \si{\newton\metre}) of the open-loop table, with time in column~1
        (\texttt{Ind\_Breakpoint}=1).
  \item A new table file \texttt{OLInput\_ROSCO.dat} was written with the
        prescribed torque time series.
\end{itemize}
The prescribed torque is a sinusoidal modulation about the Region-2 operating
torque \(T_0\),
\begin{equation}
  T_{\mathrm{gen}}(t) = T_0\,\big[\,1 + a\,
    \sin\!\big(2\pi f\,(t-t_{\mathrm{on}})\big)\,\big]
    \mathbb{1}[t\ge t_{\mathrm{on}}],\qquad a = 0.10,
  \label{eq:se-ol-forcing}
\end{equation}
i.e. a \(\pm\SI{10}{\percent}\) generator-torque oscillation. This overrides
ROSCO's torque regulation entirely, so the runs are kept short and the rotor
speed monitored for drift, with \(T_0\) matched to the aerodynamic torque at the
operating wind speed so that the net drift over the window is negligible.

\paragraph{Did it work, and what did it show.}
It worked, and it is the mechanism behind the frequency sweep of
\S\ref{sec:tower-ss-resonance}. Injecting \eqref{eq:se-ol-forcing} at the modal
frequency produced a tower-top side-to-side acceleration (\texttt{YawBrTAyp})
that \emph{builds up} to a sustained forced amplitude, rather than the free decay
seen in all the ignored \texttt{GenSpdOrTrq}/\texttt{ElecPwrCom} attempts ---
direct confirmation that the torque now reaches the structure. Sweeping \(f\)
gave a sharp resonance peaking in the SS band
(\SIrange{0.234}{0.245}{\hertz}), with a peak side-to-side acceleration of
\(\approx\SI{0.26}{\metre\per\second\squared}\) and roughly \(60\times\)
amplification over the off-resonance floor (Table~\ref{tab:ss-sweep}). The
important caveat, stated already in \S\ref{subsec:ss-torque-path}, is that the
torque here is \emph{imposed} through the open-loop table rather than produced by
a closed grid\(\to\)torque feedback: the experiment proves the \emph{structural}
frequency selectivity of the OpenFAST tower, while the reduced model
(\S\ref{subsec:se-reduced}) supplies the genuinely closed electrical drive. The
two are complementary --- the FMU shows the real tower amplifies a torque
oscillation, and the reduced model shows the grid can supply that oscillation.

\subsection{Why side-to-side and not fore-aft}
\label{subsec:se-selectivity}

The two first tower modes are nearly degenerate in frequency
(\(f_{\mathrm{ss}}\approx f_{\mathrm{fa}}\approx\SI{0.234}{\hertz}\)), so what
separates them under an electrical disturbance is not frequency but the
\emph{drive path}. The side-to-side mode is driven by the generator torque, which
carries the grid disturbance directly (\eqref{eq:se-path}). The fore-aft mode is
driven by the rotor thrust \(F_T\), which a grid disturbance can reach only
indirectly, through the drivetrain speed,
\begin{equation}
  \text{grid}\to P_e\to T_e\to \omega_m\to \lambda\to C_T\to F_T
  \to q_{\mathrm{fa}},
  \label{eq:se-fa-path}
\end{equation}
i.e. a second-order path gated by the very large rotor inertia. At a fixed wind
speed the thrust barely moves when the grid is perturbed, so the fore-aft
response to an electrical disturbance is much smaller than the side-to-side
response. This asymmetry is exactly why the SS mode --- and not the frequency-
adjacent FA mode --- is the electrically resonance-prone one.

The reduced model makes this quantitative. Because the two modes are switched
independently and share the same electrical solution
(\S\ref{subsec:stm-forward}), a single co-simulation exposes both responses to
the identical disturbance. Under a sustained load oscillation at
\SI{0.234}{\hertz} the side-to-side lock-in acceleration exceeds the fore-aft one
by a factor of \(\sim 30\) (\S\ref{subsec:nt-resonance}), and under a discrete
grid event the ratio of peak accelerations is of the same order
(\S\ref{subsec:nt-gttrip}). The tower acts as a mechanical band-pass filter on
the electrical disturbance, and the generator-torque path selects the
side-to-side mode as its output. A lock-in / contrast analysis of the two
accelerations at the forcing frequency is provided by
\texttt{casestudies/dyn\_sim/\_contrast\_ss\_fa.py}.

\subsection*{Sources}
\label{subsec:se-sources}

\begin{itemize}
  \item ROSCO open-loop control (\texttt{OL\_Mode}, \texttt{Ind\_GenTq},
        \texttt{OLInput\_ROSCO.dat}) and the internal \(K\omega^2\) torque law:
        Abbas et al.\ (2022), \emph{A reference open-source controller for fixed
        and floating offshore wind turbines}, Wind Energy Science 7, 53--73,
        \href{https://doi.org/10.5194/wes-7-53-2022}{doi.org/10.5194/wes-7-53-2022};
        documentation \href{https://rosco.readthedocs.io/}{rosco.readthedocs.io}
        (see the \emph{Open Loop Input} page); source
        \href{https://github.com/NREL/ROSCO}{github.com/NREL/ROSCO}.
  \item OpenFAST (FMU export, \texttt{YawBrTAyp} tower-top acceleration output):
        \href{https://openfast.readthedocs.io/}{openfast.readthedocs.io}.
\end{itemize}
